% Part: normal-modal-logic
% Chapter: tableaux
% Section: soundness

\documentclass[../../../include/open-logic-section]{subfiles}

\begin{document}

\olfileid[hi]{nml}{tab}{sou}

\olsection{\Log{K} के लिए यथार्थता}

\begin{editorial}
  यह यथार्थता-प्रमाण शास्त्रीय प्रतिज्ञप्तिक तर्कशास्त्र के यथार्थता-प्रमाण
  का पुनः उपयोग करता है, अर्थात् सब कुछ आरंभ से सिद्ध करता है। यदि आपको
  स्वयंपूर्ण यथार्थता-प्रमाण चाहिए, तो यह उचित है। यदि आप सामान्य टैब्लो
  के लिए यथार्थता पहले देख चुके हैं, तो यह पुनरावृत्तिपूर्ण होगा। योजना है
  कि स्वयंपूर्ण संस्करण और गैर-मोडल स्थिति पर आधारित संस्करण के बीच चयन
  संभव बनाया जाए।
\end{editorial}

\begin{explain}
  उपसर्गयुक्त !!{tableau} की यथार्थता दिखाने के लिए हमें दिखाना है कि यदि
  \[
  \sFmla{\True}{!B_1}[1], \dots, \sFmla{\True}{!B_n}[1], \sFmla{\False}{!A}[1]
  \]
  का कोई संवृत !!{tableau} है, तो $!B_1, \dots, !B_n \Entails !A$।
  प्रतिधनात्मक सिद्ध करना सरल है: यदि किसी $\mModel{M}$ और जगत~$w$ के
  लिए सभी $i=1$, \dots,~$n$ हेतु $\mSat{M}{!B_i}[w]$, किंतु
  $\mSat{M}{!A}[w]$, तो कोई !!{tableau} संवृत नहीं हो सकता। ऐसा
  प्रतिप्रतिरूप दिखाता है कि !!{tableau} की आरंभिक मान्यताएँ संतोषणीय हैं।
  प्रमाण की रणनीति यह दिखाना है कि जब भी किसी !!a{tableau} की शाखा के सभी
  उपसर्गयुक्त !!{formula} संतोषणीय हों, तब नियम का कोई भी अनुप्रयोग कम-से-कम
  एक ऐसी विस्तृत शाखा देता है जो संतोषणीय भी है। चूँकि संवृत शाखाएँ
  असंतोषणीय हैं, उपसर्गयुक्त !!{formula}ों के किसी संतोषणीय समुच्चय के हर
  !!{tableau} में कम-से-कम एक खुली शाखा होनी चाहिए।

  इस रणनीति को मोडल स्थिति में लगाने के लिए हमें ``संतोषणीय'' की अपनी
  परिभाषा को मोडल प्रतिरूपों और उपसर्गों तक विस्तृत करना होगा। ऐसा कर लेने
  पर प्रमाण सीधा है।
\end{explain}

\begin{defn}
  $P$ उपसर्गों का कोई समुच्चय हो, अर्थात् $P \subseteq (\PosInt)^*
  \setminus \{\emptyseq\}$, और $\mModel{M}$ कोई प्रतिरूप हो। कोई
  फलन~$f\colon P \to W$, $\mModel{M}$ में $P$ की \emph{व्याख्या} है यदि
  जब भी $\sigma$ और $\sigma.n$, दोनों $P$ में हों, तब
  $Rf(\sigma)f(\sigma.n)$।

  उपसर्गों~$P$ की किसी व्याख्या के सापेक्ष हम परिभाषित कर सकते हैं:
  \begin{enumerate}
  \item $\mModel{M}$, $\sFmla{\True}{!A}[\sigma]$ को तभी संतुष्ट करता है जब
    $\mSat{M}{!A}[f(\sigma)]$.
  \item $\mModel{M}$, $\sFmla{\False}{!A}[\sigma]$ को तभी संतुष्ट करता है जब
    $\mSat/{M}{!A}[f(\sigma)]$.
  \end{enumerate}
\end{defn}

\begin{defn}
  $\Gamma$ उपसर्गयुक्त !!{formula}ों का समुच्चय हो और $P(\Gamma)$ उसमें
  आने वाले उपसर्गों का समुच्चय हो। यदि $f$, $\mModel{M}$ में
  $P(\Gamma)$ की कोई व्याख्या है, तो हम कहते हैं कि $\mModel{M}$,
  $f$ के सापेक्ष $\Gamma$ को संतुष्ट करता है, $\mSat{M}{\Gamma}[f]$, यदि
  $\mModel{M}$, $f$ के सापेक्ष $\Gamma$ के प्रत्येक उपसर्गयुक्त !!{formula}
  को संतुष्ट करता हो। $\Gamma$ \emph{संतोषणीय} तभी है जब ऐसा कोई
  प्रतिरूप~$\mModel{M}$ और $P(\Gamma)$ की व्याख्या~$f$ हो कि
  $\mSat{M}{\Gamma}[f]$।
\end{defn}

\begin{prop}
  यदि $\Gamma$ में किसी !!{formula}~$!A$ और उपसर्ग~$\sigma$ के लिए
  $\sFmla{\True}{!A}[\sigma]$ और $\sFmla{\False}{!A}[\sigma]$, दोनों
  हैं, तो $\Gamma$ असंतोषणीय है।
\end{prop}

\begin{proof}
  ऐसा कोई प्रतिरूप~$\mModel{M}$ और $P(\Gamma)$ की व्याख्या~$f$ नहीं हो
  सकते कि $\mSat{M}{!A}[f(\sigma)]$ और
  $\mSat/{M}{!A}[f(\sigma)]$, दोनों हों।
\end{proof}

\begin{thm}[यथार्थता]
  \ollabel{thm:tableau-soundness}
  यदि $\Gamma$ का कोई संवृत !!{tableau} है, तो $\Gamma$ असंतोषणीय है।
\end{thm}

\begin{proof}
हम किसी !!a{tableau} की शाखा को संतोषणीय तभी कहते हैं जब उस पर स्थित
!!{signed formula}ों का समुच्चय संतोषणीय हो; और किसी !!a{tableau} को
संतोषणीय कहें यदि उसमें कम-से-कम एक संतोषणीय शाखा हो।

हम निम्न दिखाते हैं: किसी संतोषणीय !!{tableau} को अनुमान-नियमों में से किसी
एक से विस्तृत करने पर सदा संतोषणीय !!{tableau} मिलता है। इससे प्रमेय सिद्ध
होगा: कोई भी संवृत !!{tableau}, केवल $\Gamma$ की मान्यताओं वाले !!{tableau}
पर अनुमान-नियम लगाकर मिलता है। अतः यदि $\Gamma$ संतोषणीय होता, तो उसका हर
!!{tableau} संतोषणीय होता। किंतु कोई संवृत !!{tableau} स्पष्टतः संतोषणीय
नहीं है, क्योंकि उसकी सभी शाखाएँ संवृत हैं और संवृत शाखाएँ असंतोषणीय हैं।

मान लें हमारे पास कोई संतोषणीय !!{tableau} है, अर्थात् कम-से-कम एक
संतोषणीय शाखा वाला !!a{tableau}। अनुमान-नियम लगाने पर या तो किसी शाखा में
!!{signed formula} जुड़ते हैं, या शाखा दो भागों में विभक्त होती है। यदि
!!{tableau} में ऐसी संतोषणीय शाखा है जिसे विचाराधीन नियम-अनुप्रयोग विस्तृत
नहीं करता, तो वह विस्तृत !!{tableau} में संतोषणीय शाखा बनी रहती है, अतः
विस्तृत टैब्लो संतोषणीय है। इसलिए हमें केवल उस स्थिति पर विचार करना है
जहाँ नियम किसी संतोषणीय शाखा पर लगाया जाता है।

उस शाखा के !!{signed formula}ों का समुच्चय $\Gamma$ हो, और
$\sFmla{S}{!A}[\sigma] \in \Gamma$ वह !!{signed formula} हो जिस पर नियम
लगाया जाता है। यदि नियम शाखा को विभक्त नहीं करता, तो हमें दिखाना है कि
विस्तृत शाखा, अर्थात् नियम के निष्कर्षों सहित $\Gamma$, अब भी संतोषणीय
है। यदि नियम शाखा को विभक्त करता है, तो हमें दिखाना है कि उससे मिली दो
शाखाओं में कम-से-कम एक संतोषणीय है।
\tagfalse{prvDiamond}
पहले हम केवल एक आधार-वाक्य वाले संभावित अनुमानों पर विचार करते हैं।
\begin{enumerate}
\item शाखा को $\sFmla{\True}{\lnot !B}[\sigma] \in \Gamma$ पर
  $\TRule{\True}{\lnot}$ लगाकर विस्तृत किया जाता है। तब विस्तृत शाखा में
  !!{signed formula} $\Gamma \cup \{\sFmla{\False}{!B}[\sigma]\}$ हैं।
  मान लें $\mSat{M}{\Gamma}[f]$। विशेषतः
  $\mSat{M}{\lnot !B}[f(\sigma)]$। अतः
  $\mSat/{M}{!B}[f(\sigma)]$, अर्थात् $\mModel{M}$, $f$ के सापेक्ष
  $\sFmla{\False}{!B}[\sigma]$ को संतुष्ट करता है।
\item शाखा को $\sFmla{\False}{\lnot !B}[\sigma] \in \Gamma$ पर
  $\TRule{\False}{\lnot}$ लगाकर विस्तृत किया जाता है: अभ्यास।
\item शाखा को $\sFmla{\True}{!B \land !C}[\sigma] \in \Gamma$ पर
  $\TRule{\True}{\land}$ लगाकर विस्तृत किया जाता है, जिससे शाखा पर दो नए
  !!{signed formula}, $\sFmla{\True}{!B}[\sigma]$ और
  $\sFmla{\True}{!C}[\sigma]$, मिलते हैं। मान लें
  $\mSat{M}{\Gamma}[f]$, विशेषतः $\mSat{M}{!B \land !C}[f(\sigma)]$।
  तब $\mSat{M}{!B}[f(\sigma)]$ और $\mSat{M}{!C}[f(\sigma)]$। इसका अर्थ
  है कि $\mModel{M}$, $f$ के सापेक्ष $\sFmla{\True}{!B}[\sigma]$ और
  $\sFmla{\True}{!C}[\sigma]$, दोनों को संतुष्ट करता है।
\item शाखा को $\sFmla{\False}{!B \lor !C} \in \Gamma$ पर
  $\TRule{\False}{\lor}$ लगाकर विस्तृत किया जाता है: अभ्यास।
\item शाखा को $\sFmla{\False}{!B \lif !C}[\sigma] \in \Gamma$ पर
  $\TRule{\False}{\lif}$ लगाकर विस्तृत किया जाता है। इससे शाखा पर दो नए
  !!{signed formula}, $\sFmla{\True}{!B}[\sigma]$ और
  $\sFmla{\False}{!C}[\sigma]$, मिलते हैं। मान लें
  $\mSat{M}{\Gamma}[f]$, विशेषतः
  $\mSat/{M}{!B \lif !C}[f(\sigma)]$। तब
  $\mSat{M}{!B}[f(\sigma)]$ और $\mSat/{M}{!C}[f(\sigma)]$। इसका अर्थ
  है कि $\mModel{M}, f$, $\sFmla{\True}{!B}[\sigma]$ और
  $\sFmla{\False}{!C}[\sigma]$, दोनों को संतुष्ट करते हैं।

\iftag{prvBox}{%
\item शाखा को $\sFmla{\True}{\Box !B}[\sigma] \in \Gamma$ पर
  $\TRule{\True}{\Box}$ लगाकर विस्तृत किया जाता है:
  \iftag{probBox}{अभ्यास।}{इससे शाखा पर किसी $\sigma.n \in P(\Gamma)$
    के लिए नया !!{signed formula}~$\sFmla{\True}{!B}[\sigma.n]$ मिलता है
    (क्योंकि $\sigma.n$ प्रयुक्त होना चाहिए)। मान लें
    $\mSat{M}{\Gamma}[f]$, विशेषतः
    $\mSat{M}{\Box !B}[f(\sigma)]$। चूँकि $f$ उपसर्गों की व्याख्या है और
    $\sigma$, $\sigma.n$, दोनों $P(\Gamma)$ में हैं, हम जानते हैं कि
    $Rf(\sigma)f(\sigma.n)$। अतः $\mSat{M}{!B}[f(\sigma.n)]$, अर्थात्
    $\mModel{M}, f$, $\sFmla{\True}{!B}[\sigma.n]$ को संतुष्ट करते हैं।}

\item शाखा को $\sFmla{\False}{\Box !B}[\sigma] \in \Gamma$ पर
  $\TRule{\False}{\Box}$ लगाकर विस्तृत किया जाता है:
  \iftag{probBox}{अभ्यास।}{इससे नया !!{signed formula}~$\sFmla{\False}{!A}[\sigma.n]$ मिलता है, जहाँ $\sigma.n$
    शाखा पर नया उपसर्ग है, अर्थात् $\sigma.n \notin P(\Gamma)$। चूँकि
    $\Gamma$ संतोषणीय है, ऐसी कोई $\Struct{M}$ और $P(\Gamma)$ की
    व्याख्या~$f$ है कि $\Sat{M}{\Gamma}[f]$, विशेषतः
    $\mSat/{M}{\Box !B}[f(\sigma)]$। हमें दिखाना है कि $\Gamma \cup
    \{\sFmla{\False}{!B}[\sigma.n]\}$ संतोषणीय है। इसके लिए हम
    $P(\Gamma) \cup \{\sigma.n\}$ की व्याख्या इस प्रकार परिभाषित करते हैं:

  चूँकि $\mSat/{M}{\Box !B}[f(\sigma)]$, ऐसा $w \in W$ है कि
  $Rf(\sigma)w$ और $\mSat/{M}{!B}[w]$। $f'$ को $f$ जैसा रखें, बस
  $f'(\sigma.n) = w$। चूँकि $f'(\sigma) = f(\sigma)$ और
  $Rf(\sigma)w$, इसलिए $Rf'(\sigma)f'(\sigma.n)$; अतः $f'$,
  $P(\Gamma) \cup \{\sigma.n\}$ की व्याख्या है। स्पष्टतः
  $\mSat/{M}{!B}[f'(\sigma.n)]$। सभी उपसर्गों $\sigma' \in P(\Gamma)$
  के लिए $f(\sigma') = f'(\sigma')$ होने से
  $\mSat{M}{\Gamma}[f']$। अतः $\mModel{M}, f'$, $\Gamma \cup
  \{\sFmla{\False}{!B}[\sigma.n]\}$ को संतुष्ट करते हैं।}
  }{}

\iftag{prvDiamond}{%
\item शाखा को $\sFmla{\True}{\Diamond !B}[\sigma] \in \Gamma$ पर
  $\TRule{\True}{\Diamond}$ लगाकर विस्तृत किया जाता है:
  \iftag{probDiamond}{अभ्यास।}{इससे नया !!{signed formula}~$\sFmla{\True}{!A}[\sigma.n]$ मिलता है, जहाँ $\sigma.n$
    शाखा पर नया उपसर्ग है, अर्थात् $\sigma.n \notin P(\Gamma)$। चूँकि
    $\Gamma$ संतोषणीय है, ऐसी कोई $\Struct{M}$ और $P(\Gamma)$ की
    व्याख्या~$f$ है कि $\Sat{M}{\Gamma}[f]$, विशेषतः
    $\mSat{M}{\Diamond !B}[f(\sigma)]$। हमें दिखाना है कि $\Gamma \cup
    \{\sFmla{\True}{!B}[\sigma.n]\}$ संतोषणीय है। इसके लिए हम
    $P(\Gamma) \cup \{\sigma.n\}$ की व्याख्या इस प्रकार परिभाषित करते हैं:

  चूँकि $\mSat{M}{\Diamond !B}[f(\sigma)]$, ऐसा $w \in W$ है कि
  $Rf(\sigma)w$ और $\mSat{M}{!B}[w]$। $f'$ को $f$ जैसा रखें, बस
  $f'(\sigma.n) = w$। चूँकि $f'(\sigma) = f(\sigma)$ और
  $Rf(\sigma)w$, इसलिए $Rf'(\sigma)f'(\sigma.n)$; अतः $f'$,
  $P(\Gamma) \cup \{\sigma.n\}$ की व्याख्या है। स्पष्टतः
  $\mSat{M}{!B}[f'(\sigma.n)]$। सभी उपसर्गों $\sigma' \in P(\Gamma)$
  के लिए $f(\sigma') = f'(\sigma')$ होने से
  $\mSat{M}{\Gamma}[f']$। अतः $\mModel{M}, f'$, $\Gamma \cup
  \{\sFmla{\True}{!B}[\sigma.n]\}$ को संतुष्ट करते हैं।}
\item शाखा को $\sFmla{\False}{\Diamond !B}[\sigma] \in \Gamma$ पर
  $\TRule{\False}{\Diamond}$ लगाकर विस्तृत किया जाता है:
  \iftag{probDiamond}{अभ्यास।}{इससे शाखा पर किसी $\sigma.n \in
    P(\Gamma)$ के लिए नया !!{signed formula}~$\sFmla{\False}{!B}[\sigma.n]$ मिलता है (क्योंकि $\sigma.n$
    प्रयुक्त होना चाहिए)। मान लें $\mSat{M}{\Gamma}[f]$, विशेषतः
    $\mSat/{M}{\Diamond !B}[f(\sigma)]$। चूँकि $f$ उपसर्गों की व्याख्या
    है और $\sigma$, $\sigma.n$, दोनों $P(\Gamma)$ में हैं, हम जानते हैं
    कि $Rf(\sigma)f(\sigma.n)$। अतः $\mSat/{M}{!B}[f(\sigma.n)]$,
    अर्थात् $\mModel{M}, f$, $\sFmla{\False}{!B}[\sigma.n]$ को संतुष्ट
    करते हैं।}
  }{}
\end{enumerate}
अब दो आधार-वाक्यों वाले संभावित अनुमानों पर विचार करें।
\begin{enumerate}
\item शाखा को $\sFmla{\False}{!B \land !C}[\sigma] \in \Gamma$ पर
  $\TRule{\False}{\land}$ लगाकर विस्तृत किया जाता है। इससे दो शाखाएँ
  मिलती हैं: $\sFmla{\False}{!B}[\sigma]$ से आगे जाने वाली बाईं शाखा और
  $\sFmla{\False}{!C}[\sigma]$ से आगे जाने वाली दाईं शाखा। मान लें
  $\mSat{M}{\Gamma}[f]$, विशेषतः
  $\mSat/{M}{!B \land !C}[f(\sigma)]$। तब
  $\mSat/{M}{!B}[f(\sigma)]$ या $\mSat/{M}{!C}[f(\sigma)]$। पहली स्थिति
  में $\mModel{M}, f$, $\sFmla{\False}{!B}[\sigma]$ को संतुष्ट करते हैं,
  अर्थात् बाईं शाखा संतोषणीय है। दूसरी स्थिति में $\mModel{M}, f$,
  $\sFmla{\False}{!C}[\sigma]$ को संतुष्ट करते हैं, अर्थात् दाईं शाखा
  संतोषणीय है।
\item शाखा को $\sFmla{\True}{!B \lor !C}[\sigma] \in \Gamma$ पर
    $\TRule{\True}{\lor}$ लगाकर विस्तृत किया जाता है: अभ्यास।
\item शाखा को $\sFmla{\True}{!B \lif !C}[\sigma] \in \Gamma$ पर
    $\TRule{\True}{\lif}$ लगाकर विस्तृत किया जाता है: अभ्यास।
\end{enumerate}
\end{proof}

\begin{prob}
\olref[nml][tab][sou]{thm:tableau-soundness} का प्रमाण पूरा कीजिए।
\end{prob}

\begin{cor}
\ollabel{cor:entailment-soundness}
यदि $\Gamma \Proves !A$, तो $\Gamma \Entails !A$।
\end{cor}

\begin{proof}
  यदि $\Gamma \Proves !A$, तो किन्हीं $!B_1$, \dots, $!B_n \in
  \Gamma$ के लिए $\Delta = \{\sFmla{\False}{!A}[1],
  \sFmla{\True}{!B_1}[1], \dots, \sFmla{\True}{!B_n}[1]\}$ का कोई
  संवृत !!{tableau} है। हम दिखाना चाहते हैं कि $\Gamma \Entails !A$। मान
  लें ऐसा नहीं है; तब किसी $\mModel{M}$ और $w$ के लिए $i=1$, \dots,~$n$
  हेतु $\mSat{M}{!B_i}[w]$, किंतु $\mSat/{M}{!A}[w]$। $f(1) = w$ रखें;
  तब $f$, $\mModel{M}$ में $P(\Delta)$ की व्याख्या है, और $\mModel{M}$,
  $f$ के सापेक्ष $\Delta$ को संतुष्ट करता है। किंतु
  \olref{thm:tableau-soundness} से $\Delta$ असंतोषणीय है, क्योंकि उसका
  संवृत !!{tableau} है; यह विरोधाभास है। अतः अंततः $\Gamma \Proves !A$
  होना ही चाहिए।
\end{proof}

\begin{cor}
\ollabel{cor:weak-soundness}
यदि $\Proves !A$, तो $!A$ सभी प्रतिरूपों में सत्य है।
\end{cor}

\end{document}
